\chapter*{Comment lire ce document}
\addcontentsline{toc}{chapter}{Comment lire ce document}
\markboth{COMMENT LIRE CE DOCUMENT}{}

\section*{Le principe : un empilement}

On ne présente pas les outils par ordre alphabétique ni par importance. On les
pose dans l'ordre où il a fallu les poser, et chaque couche ne peut exister que
parce que la précédente est là.

\begin{figure}[htbp]
\centering
\resizebox{\textwidth}{!}{%
\begin{tikzpicture}[font=\small,
  c/.style={draw=grisligne,fill=grisfond,rounded corners=3pt,minimum width=11cm,
            minimum height=9mm,align=center},
  n/.style={font=\scriptsize\bfseries,insared}]
\foreach \y/\t/\o in {
  0/{Le socle}/{Git · Python · venv},
  1/{L'application web}/{Flask · Jinja2 · Werkzeug · Gunicorn · nginx · Tailwind},
  2/{Les données}/{SQLAlchemy · Alembic · SQLite · PostgreSQL · Fernet},
  3/{Les machines louées}/{Docker · supervisord · OpenSSH · ttyd · Xvnc · XFCE · noVNC},
  4/{Le pilotage du parc}/{subprocess · étiquettes · agent · Faucheur · systemd},
  5/{L'infrastructure}/{Vagrant · VMware · Ansible · Vault · inventaire dynamique},
  6/{La qualité}/{pytest · ruff · bandit · Trivy · Syft · GitHub Actions}} {
  \node[c] at (0,\y*1.15) {\textbf{\t}\\[1pt]\textcolor{gristexte}{\scriptsize\o}};
  \node[n] at (-6.3,\y*1.15) {\y};
}
\draw[-{Stealth[length=3mm]},insared,thick] (6.3,-0.45) -- (6.3,7.4)
  node[midway,right=4pt,insared,font=\scriptsize,align=left]
  {chaque couche\\suppose\\la précédente};
\end{tikzpicture}}
\caption{Les sept couches. On ne peut pas conteneuriser avant d'avoir une
application, ni automatiser le déploiement avant d'avoir quelque chose à
déployer.}
\end{figure}

\section*{Chaque outil, quatre questions}

Toutes les fiches ont la même forme. Vous pouvez lire le document en entier,
ou n'en consulter qu'une fiche.

\begin{description}
\item[C'est quoi] en une ou deux phrases, sans supposer qu'on connaît.
\item[Pourquoi celui-là] quelles étaient les autres possibilités, et ce qui a
  fait pencher. C'est la partie la plus utile : un outil se comprend par ce
  qu'on a écarté.
\item[Son rôle ici] précisément, dans ce projet, avec le fichier concerné.
\item[Ce que ça a changé] ce que la plateforme savait faire avant, et après.
\end{description}

\begin{retenir}
Si vous ne devez retenir qu'une chose par outil, retenez le \textbf{pourquoi
celui-là}. Savoir qu'on utilise Gunicorn ne vaut rien ; savoir pourquoi le
serveur de développement de Flask ne suffit pas, si.
\end{retenir}

\section*{Ce que la plateforme doit savoir faire}

Pour mesurer l'apport de chaque brique, il faut savoir où l'on va. La
plateforme finale doit :

\begin{enumerate}
\item accueillir des comptes utilisateurs et les authentifier ;
\item proposer un catalogue de systèmes Linux ;
\item créer une machine à la demande, en quinze secondes ;
\item en donner l'accès par SSH, par un terminal web et par un bureau
  graphique ;
\item la détruire toute seule à l'échéance ;
\item répartir les machines sur plusieurs serveurs ;
\item survivre à la panne de l'un d'eux ;
\item et se redéployer entièrement en une commande.
\end{enumerate}

Aucune de ces huit capacités n'est apportée par un seul outil. Chacune naît de
la rencontre de plusieurs briques --- et c'est précisément ce que ce document
cherche à rendre lisible.
